The experimental results demonstrate the effectiveness of using symbolic computational tools to analyze fundamental signal operations. By using the MATLAB Symbolic Math Toolbox, the mathematical construction of signals $x(t)$ and $y(t)$ maintained infinite precision, successfully avoiding the discretization errors typically found in numerical methods. The analysis of cross-correlation $R_{xy}(\tau)$ specifically highlighted the critical role of boundary conditions in piecewise integration. As visualized in the transition from partial to full overlap in Cases 2 through 4, the resulting symmetrical waveform provides a precise quantification of signal similarity over the time-lag $\tau$. This aligns with established theory that cross-correlation serves as a measure of the degree to which two signals are correlated as one is shifted relative to the other \cite{oppenheim1997}.

Furthermore, the study of Linear Time-Invariant (LTI) system responses through convolution revealed the distinct characteristics of causal system behavior. The interaction between the input pulse $x_{in}(t)$ and the exponential impulse response $h(t) = e^{-10t}u(t)$ resulted in an output $y_{out}(t)$ that accurately reflected the system's time constant. The sharp rise and subsequent rapid nonlinear decay confirm the system's immediate and fading memory response to an external stimulus \cite{haykin1999}. Such behavior is a hallmark of first-order LTI systems, where the output is essentially a weighted integral of past inputs, governed by the specific decay rate of the impulse response \cite{proakis2007}.

In conclusion, this laboratory exercise successfully verified the theoretical underpinnings of signal interaction and system characterization. The structured methodology, moving from mathematical modeling to segmented computational execution, ensured that the resulting data was both mathematically grounded and highly accurate. The ability to visualize these complex interactions, such as the stages of signal passing and the processing of a pulse by an exponential decay, reinforces the importance of using computational logic to solve calculus-heavy signal processing tasks. These findings serve as a practical validation of the core principles of signals and systems, providing a reliable framework for analyzing more complex waveform interactions in future engineering applications.